\documentclass{article}
\usepackage[T1]{fontenc}
\usepackage{lmodern}
\usepackage{graphicx}
\usepackage{amsmath,amssymb}
\usepackage[a4paper,margin=2cm]{geometry}
\usepackage[absolute,overlay]{textpos}
\setlength{\TPHorizModule}{1mm}
\setlength{\TPVertModule}{1mm}
\usepackage[indonesian]{babel}
\usepackage{hyperref}
\setlength{\emergencystretch}{2em}
% unit: Halaman 28

\begin{document}

% === Halaman 28 (python/native) ===

Selanjutnya, XOR-kan semua hasil antara tersebut:


(02 • 26)   = 4C =  0100 1100


(03 • 7B)   = 8D = 1000 1101

       (01 • BD)  = BD = 1011 1101

       (01 • 43)   = 43 =  0100 0011 


                               0011 1111 = 3F

Jadi, (02 • 26)  (03 • 7B)  (01 • BD)  (01 • 43) = 3F

Persamaan lainnya diselesaikan dengan cara yang sama.

(02 • 26)  (03 • 7B)  (01 • BD)  (01 • 43) = 3F

\end{document}
